% Chapter 3 - Methodology.
% Sources: docs/plan.md, docs/environment.md, claude/desarrollo.md,
% claude/diario.md, .github/workflows/*.yml.
% Target length: 6-8 pages.

\chapter{Methodology}
\label{ch:metodologia}

This chapter describes the working method behind the proof of
concept. It is not a generic overview of agile or test-driven
development: the goal is to record the specific choices that were
made and the reasons for them, so that any reader can reproduce the
process or take it as a basis for future iterations.

The method rests on four pillars. First, a phased process with a
formal deliverable per phase, complemented by short iterations
(sprints) inside the implementation phase. Second, a single
reproducible development environment that works identically on a
local machine, on a remote server and on cloud workspaces. Third, a
strict code-quality discipline supported by continuous integration.
Fourth, an explicit and traceable use of AI-assisted development.

The four pillars are tied together by a set of versioned context
files that act as the operational memory of the project: the
chronological journal, the technical-decisions log, the pending-ideas
backlog and the development guidelines. These files live in the
\texttt{claude/} directory of the repository and are loaded by the
AI assistant at the start of every session, but they also serve as
the only written record of the rationale behind every design and
code change.

%==================================================================
\section{Process framework: phases and sprints}
\label{sec:proceso}

The work is organized in seven phases, each with a specific
deliverable. \Cref{tab:phases} lists them. The structure follows the
classical engineering breakdown (problem definition, requirements,
architecture, technology selection, environment, implementation,
evaluation) and assigns each phase to a stable artefact in the
repository. The first five phases are sequential because their
outputs feed the next; the sixth phase contains internal iterations,
and the seventh runs in parallel with the sixth from the moment the
documentation work begins.

\begin{table}[ht]
  \centering
  \small
  \begin{tabular}{p{0.07\linewidth} p{0.30\linewidth} p{0.30\linewidth} p{0.20\linewidth}}
    \toprule
    \textbf{Phase} & \textbf{Activity} & \textbf{Deliverable} & \textbf{Status} \\
    \midrule
    1 & Problem definition          & \texttt{docs/vision.md}        & Completed \\
    2 & Requirements analysis       & \texttt{docs/requirements.md}  & Completed \\
    3 & Architecture design         & \texttt{docs/architecture.md} + \texttt{docs/diagrams.md} & Completed \\
    4 & Technology selection        & \texttt{docs/technologies.md}  & Completed \\
    5 & Environment configuration   & Working devcontainer + \texttt{docs/environment.md} & Completed \\
    6 & Iterative development (TDD) & Working software (sprints 1--3) & In progress \\
    7 & Evaluation and report       & This document                   & In progress \\
    \bottomrule
  \end{tabular}
  \caption{The seven phases of the project and their deliverables.}
  \label{tab:phases}
\end{table}

Inside Phase 6, the work is split into \emph{sprints} of two to
three weeks each, with a clear scope agreed at the start and a
formal close in the journal. Three sprints were completed during the
development period:

\begin{itemize}
  \item \textbf{Sprint 1} (Minimum viable core, CLI). Delivered the
        descriptor tree, the cascade-inheritance model, the
        connector contract, the mock connector, the YAML loader, the
        Job engine and a working CLI. Closed with a runnable demo
        for the supervisor.
  \item \textbf{Sprint 2} (Resolver and template closure). Delivered
        the lexical-closure resolver, the connector type registry
        and the live-session pool. Closed with the demo tree
        deployable end-to-end without any inline hypervisor
        information in the descriptors.
  \item \textbf{Sprint 3} (REST daemon and packaged distribution).
        Delivered the FastAPI daemon, the pure-HTTP CLI client, the
        installable Docker image hosted in the GitHub Container
        Registry, a single-line installer script and a tutorial for
        the supervisor.
\end{itemize}

The journal (\texttt{claude/diario.md}) records every decision in
chronological order; the technical-decisions log
(\texttt{claude/decisiones.md}) keeps the numbered DEC entries; both
are cross-referenced. The proof of concept reaches the end of
Sprint~3 with 244 automated tests passing, 2{,}522 lines of source
code under \texttt{src/univorch/} and 2{,}796 lines of tests under
\texttt{tests/}. \Cref{ch:desarrollo} describes the delivered
software in detail.

%==================================================================
\section{Reproducible development environment}
\label{sec:entorno}

The development environment is defined by a single
\texttt{.devcontainer/devcontainer.json} file in the repository. The
same file produces an identical environment in three setups:

\begin{itemize}
  \item Visual Studio Code with local Docker (the standard option).
  \item Visual Studio Code with the \emph{Remote~--~SSH} extension
        connected to a Linux server, with Docker running on the
        server. This is the recommended option for Phase 6 because
        the server has direct network access to the teaching VMware
        and Proxmox hypervisors and lifts the hour limits of cloud
        workspaces.
  \item GitHub Codespaces in the browser. Useful for short sessions
        or for working from a machine without Docker.
\end{itemize}

The base image of the devcontainer is the official Microsoft Python
3.12 image. The container installs \texttt{uv}~\cite{uv-docs} as the
package manager and synchronizes all development dependencies
declared in \texttt{pyproject.toml} with the \texttt{dev} extra.
Editor configuration, linter settings and useful VSCode extensions
are also part of the file. A developer who clones the repository and
opens it in VSCode is ready to work within minutes, without any
manual installation.

The development image and the production image are deliberately
different. The development image is convenience-oriented (full
toolchain, debugger, Claude Code extension). The production image is
a multi-stage minimal build (Python 3.12 slim, the application code
and \texttt{uv} from a layer cache) of around 200 megabytes. The
production image is published in the GitHub Container Registry
(\texttt{ghcr.io/clamaveruma/univorch}) on every tagged release.

The repository follows a Python source layout with a single
distribution package, \texttt{univorch}, and the test suite split in
unit and integration tests:

\begin{verbatim}
univorch/
  src/univorch/        application code
  tests/unit/          isolated unit tests
  tests/integration/   tests with TinyDB + MockConnector
  demo/                YAML examples and demo guide
  docs/                phase deliverables (English)
  claude/              project memory and AI context (Spanish)
  .devcontainer/       reproducible dev environment
  .github/workflows/   CI/CD pipelines
  Dockerfile           production image
  docker-compose.yml   service composition
  pyproject.toml       unified project configuration
\end{verbatim}

%==================================================================
\section{Quality discipline}
\label{sec:calidad}

The codebase was developed under a strict test-driven discipline.
Every public function has at least one test, every Command is tested
with both an accepting and a rejecting case, and every Repository is
tested against an in-memory TinyDB. The pure Resolver, being a
function over plain data, is also exercised by property-based tests
written with Hypothesis~\cite{hypothesis-docs}, which automatically
generates trees of folders and descriptors and checks invariants
such as ``a child folder cannot shadow a hypervisor that was never
imported.''

The four quality gates run on every commit. They are also enforced
in continuous integration; a failed check blocks the build.

\begin{description}
  \item[Ruff~\cite{ruff-docs}] -- linter and formatter in a single
        tool. Replaces \texttt{flake8}, \texttt{isort} and
        \texttt{black}. Configured in strict mode for the application
        source code.
  \item[mypy --strict] -- static type checker. Validates abstract
        base classes (used for the connector contract), Protocol
        types (used for the orchestrator API) and Pydantic models.
        The strict mode forces explicit annotations on every public
        function.
  \item[pytest + pytest-cov] -- test runner with branch coverage.
        244 tests pass at the close of Sprint~3. Coverage exceeds
        95\% in the central modules (resolver, service, commands,
        repositories, connectors).
  \item[Hypothesis] -- property-based testing on the Resolver and
        the inheritance resolution logic.
\end{description}

\Cref{tab:test-distribution} summarizes the distribution of tests by
type and module.

\begin{table}[ht]
  \centering
  \small
  \begin{tabular}{l r r}
    \toprule
    \textbf{Area} & \textbf{Tests} & \textbf{Style} \\
    \midrule
    Models (Pydantic)           & 28 & Unit \\
    Resolver                    & 26 & Unit + property-based \\
    Connectors (ABC + mock)     & 39 & Unit \\
    Commands                    & 35 & Unit + integration \\
    Job engine                  & 12 & Integration \\
    Repositories                & 22 & Integration (in-memory TinyDB) \\
    Service (facade)            & 39 & Integration \\
    CLI                         & 30 & Unit + integration \\
    REST daemon                 & 13 & Integration (ASGI test client) \\
    \midrule
    \textbf{Total}              & \textbf{244} & \\
    \bottomrule
  \end{tabular}
  \caption{Distribution of automated tests at the close of Sprint~3.}
  \label{tab:test-distribution}
\end{table}

The same checks run locally and in CI. The local invocation is a
single command (\texttt{uv run pytest} for tests; \texttt{uv run
ruff check}, \texttt{uv run ruff format} and \texttt{uv run mypy} for
the others) and the CI runs exactly the same commands. There is no
divergence between the developer's machine and the CI environment;
the only difference is that CI runs in a clean container at every
push, eliminating ``it works on my machine'' problems.

%==================================================================
\section{Continuous integration and delivery}
\label{sec:cicd}

Three GitHub Actions workflows automate the build and release
processes. They are stored in \texttt{.github/workflows/} and are
the single source of truth for what ``green'' means.

\textbf{\texttt{ci.yml}} runs on every push and every pull request
to \texttt{main}. It executes the four quality gates of
\Cref{sec:calidad}: Ruff lint, Ruff format check, mypy strict and
pytest. A failed gate is reported on GitHub and prevents the change
from being marked as passing.

\textbf{\texttt{publish.yml}} is triggered by an annotated Git tag
of the form \texttt{vX.Y.Z}. It builds the multi-stage production
image and publishes it to the GitHub Container Registry under two
tags: an immutable \texttt{vX.Y.Z} (which never changes) and a
mutable \texttt{latest} (which always points to the most recent
release). The choice to use signed Git tags makes the publication
explicit: only the maintainer creates a tag, and only that act
produces a new image. While the API surface is unstable, the project
stays in the \texttt{0.x} series.

\textbf{\texttt{memoria.yml}} compiles the LaTeX sources of this
report on every push to \texttt{main} that touches the
\texttt{docs/memoria/} directory. It renders the Mermaid diagrams to
PDF, compiles the document with XeLaTeX, uploads the PDF as an
artefact and pushes a copy to a dedicated \texttt{pdf-preview}
branch so that the supervisor always has a stable URL to the latest
version.

The end result is that the supervisor or any external reviewer can
verify the state of the project at any moment without contacting
the author: the public image is downloadable from \texttt{ghcr.io},
the source code is on the \texttt{main} branch of GitHub, the
quality gates are visible in the Actions tab, and the report
preview is one click away.

%==================================================================
\section{AI-assisted development}
\label{sec:ia}

This work was developed with the assistance of Claude (Anthropic),
used through both the Claude Code command-line tool and the VSCode
extension. We discuss this aspect explicitly because the use of an
AI assistant influences how the proof of concept was built and
because we believe a TFG that uses such a tool should describe how
it does so.

The model and effort level vary with the task. Quick, mechanical
tasks (commits, file renames, documentation transcription) use a
small and fast model. Routine coding work uses a mid-range model.
Critical architectural decisions are taken with the most capable
available model and the maximum effort setting. The mapping between
phase and model is recorded in the project's \texttt{CLAUDE.md}
file and is part of the methodology.

The collaboration follows a fixed pattern: \emph{propose, clarify,
confirm, execute}. The assistant proposes an approach in plain
prose, the author asks clarifying questions or pushes back when the
proposal is wrong, the author confirms the path forward, and only
then the assistant writes code or modifies files. The author keeps
final authority on every decision and is responsible for defending
this report; the assistant acts as a structured collaborator, not as
a black box. The journal records the entire reasoning of every
non-trivial decision exactly as it was discussed, so that any
section of the report can be traced back to its origin.

The assistant has access to a persistent context made of versioned
files. The entry point is \texttt{CLAUDE.md} in the repository root,
which imports four files:

\begin{itemize}
  \item \texttt{claude/proyecto.md} -- one-page description of the
        project, scope and constraints.
  \item \texttt{claude/diario.md} -- chronological journal of every
        decision in its context.
  \item \texttt{claude/decisiones.md} -- numbered log of technical
        decisions (DEC-001 to DEC-036 at the time of writing), each
        with a back-reference to the journal entry where it was
        taken.
  \item \texttt{claude/desarrollo.md} -- development guidelines:
        sprints, coding conventions, exception style, docstring
        format, working agreement with the author.
  \item \texttt{claude/ideas-pendientes.md} -- backlog of ideas
        outside the current scope, classified by area (web GUI,
        teaching layer, report material, UX).
\end{itemize}

The four files use Markdown and natural prose so they are equally
usable by a human reader: a maintainer who picks up the project
after this TFG ends can read \texttt{diario.md} as a project history
and \texttt{decisiones.md} as a design rationale, without needing
the AI assistant at all.

%==================================================================
\section{Traceability and knowledge management}
\label{sec:trazabilidad}

The set of \texttt{claude/} files is also the traceability backbone
of the project. Every numbered design decision in
\texttt{decisiones.md} points to the journal entry where it was
discussed; every journal entry references the decisions taken
during that session. The result is that any architectural or
implementation choice in the codebase can be traced back to its
motivation in a few clicks.

A controlled vocabulary is maintained in \texttt{docs/glossary.md}.
This file is the single source of truth for the terminology used in
the report, in the technical documentation and in the future
administrator manual. The glossary distinguishes carefully between
terms that look interchangeable but are not (e.g.\ ``descriptor''
and ``virtual machine'' belong to two different layers of the
vocabulary; the orchestrator manipulates descriptors, the user
sees virtual machines), and records the language convention chosen
for each: external vocabulary in user-facing components, internal
vocabulary in the code and the technical documents. The glossary of
this report is generated from the same source.

The pending-ideas file (\texttt{claude/ideas-pendientes.md}) keeps
proposals that are out of scope for the current sprint. It is
revised at the start of every sprint to decide what enters next and
what can be promoted to the report's future-work section. This
prevents valuable ideas from being lost between sprints without
mixing them with the active scope.

Three external reviewers (a hypothetical maintainer, the supervisor
and the tribunal) would all benefit from the same traceability:
they can independently audit how a decision was reached without
having to ask the author. We consider this an explicit deliverable
of the methodology, not a side effect.
